% Created 2016-11-30 Wed 10:39
\documentclass[bigger]{beamer}
\usepackage[utf8]{inputenc}
\usepackage[T1]{fontenc}
\usepackage{fixltx2e}
\usepackage{graphicx}
\usepackage{grffile}
\usepackage{longtable}
\usepackage{wrapfig}
\usepackage{rotating}
\usepackage[normalem]{ulem}
\usepackage{amsmath}
\usepackage{textcomp}
\usepackage{amssymb}
\usepackage{capt-of}
\usepackage{hyperref}
\mode<beamer>{\useinnertheme{rounded}\usecolortheme{rose}\usecolortheme{dolphin}\setbeamertemplate{navigation symbols}{}\setbeamertemplate{footline}[frame number]{}}
\mode<beamer>{\usepackage{amsmath}\usepackage{ae,aecompl} \usepackage{sgamevar,url}}
\usetheme{default}
\author{Christoph Schottmüller}
\date{}
\title{Game Theory: Hard evidence\footnote{License: CC Attribution ShareAlike 4.0}}
\hypersetup{
 pdfauthor={Christoph Schottmüller},
 pdftitle={Game Theory: Hard evidence\footnote{License: CC Attribution ShareAlike 4.0}},
 pdfkeywords={},
 pdfsubject={},
 pdfcreator={Emacs 24.5.1 (Org mode 8.3.3)}, 
 pdflang={English}}
\begin{document}

\maketitle
\begin{frame}{Outline}
\tableofcontents
\end{frame}



\section{Introduction}
\label{sec:orgheadline2}

\begin{frame}[label={sec:orgheadline1}]{Introduction}
\begin{itemize}
\item hard evidence
\begin{itemize}
\item A has private information but can disclose it truthfully to B
\end{itemize}
\item examples
\begin{itemize}
\item financial reports by companies
\item disclose income by showing tax statement
\item certify quality by an expert/rating agency (possibly costly)
\end{itemize}
\item questions
\begin{itemize}
\item what information should you (not) disclose?
\item when are mandatory disclosure or privacy rules optimal?
\end{itemize}
\end{itemize}
\end{frame}

\section{Unraveling}
\label{sec:orgheadline6}

\begin{frame}[label={sec:orgheadline3}]{A simple example}
\begin{enumerate}
\item a manager A privately learns the value of the firm \(v\in V=\{v_1,\dots,v_n\}\) with \(v_1<v_2<\dots<v_n\)
\item A can disclose his information to the board B
\item B pays a bonus to A depending on the expected value of the firm (the higher B expects \(v\) to be the higher the bonus)
\end{enumerate}

\vspace*{1cm}

What is the (perfect Bayesian) equilibrium of the game?
\end{frame}


\begin{frame}[label={sec:orgheadline4}]{A more general setup}
\begin{itemize}
\item A's payoff is \(\pi(v,b)\) where \(b\) is B's belief (probability distribution over  \(V=\{v_1,\dots,v_n\}\)) in stage 3
\end{itemize}
\begin{block}{Assumptions}
\begin{itemize}
\item \(\pi\) increases in \(v\)
\item \(\pi\) "increases" in \(b\): \(\pi(v,b)>\pi(v,b')\) if \(b\neq b'\) first order stochastically dominates \(b'\):

\(\sum_{v_j\leq v}b(v_j)\leq \sum_{v_j\leq v}b'(v_j)\qquad\text{ for all }v\)
\end{itemize}
\end{block}

\begin{itemize}
\item A can certify some subset of \(V\) in stage 2 (the true \(v\) must be in the set he certifies)
\begin{itemize}
\item example:  A can certify "\(v\) is in \(\{v_2,v_4\}\)" if \(v=v_2\)
\end{itemize}
\end{itemize}
\begin{block}{Assumption}
A can certify that the value is at least \(v\), i.e. if \(v=v_i\) then the set \(\{v_i,v_{i+1},\dots,v_n\}\) is certifiable
\end{block}
\end{frame}

\begin{frame}[label={sec:orgheadline5}]{Unraveling theorem in the more general setup}
\begin{theorem}[Unraveling]
The equilibrium involves full information revelation and pessimistic beliefs.
\end{theorem}
\begin{itemize}
\item pessimistic beliefs: for any certified set, B beliefs that \(v\) is the lowest \(v_i\) in the set
\end{itemize}

\textbf{Proof. } 
\end{frame}

\section{Partial Unraveling}
\label{sec:orgheadline13}

\begin{frame}[label={sec:orgheadline7}]{Costly verification I}
\begin{itemize}
\item say disclosure comes at cost \(c>0\)
\item for which \(v_i\) will A disclose, for which not?
\end{itemize}
\end{frame}

\begin{frame}[label={sec:orgheadline8}]{Costly verification II}
\begin{itemize}
\item let \(v'\) be the highest \(\tilde{v}\) that satisfies
\end{itemize}
\(\hspace*{2cm}\mathbb{E}[v| v\leq \tilde{v}] <c\)
\begin{itemize}
\item equilibrium:
\begin{itemize}
\item not disclose for \(v\leq v'\)
\item disclose for \(v> v'\)
\end{itemize}
\end{itemize}
\end{frame}

\begin{frame}[label={sec:orgheadline9}]{Dye's model I}
\begin{itemize}
\item back to simple disclosure: A cannot certify sets but at best the actual \(v\)
\item A knows \(v\) with probability \(q\)
\item A does not know \(v\) with probability \(1-q\)
\item prior belief over \(\{v_1,\dots,v_n\}\) is denoted by \(f\) (with full support)
\item A's payoff is \(\alpha v+\bar{v}\) where \(\alpha\geq 0\) and \(\bar v\) is B's expectation over \(v\) in stage 3
\item same timing as in simple example (note that A cannot disclose anything if he does not know \(v\))
\end{itemize}
\end{frame}

\begin{frame}[label={sec:orgheadline10}]{Dye's model II}
\begin{example}[partial disclosure]
\begin{itemize}
\item \(V=\{1,2,9\}\) and \(f(v_i)=1/3\) for \(v_i\in V\)
\item \(q=3/4\) and \(\alpha=1\)
\item Equilibrium: A discloses if \(v=9\) and does not disclose otherwise:
\begin{itemize}
\item given this, what is the prob of non-disclosure?

\item if non-disclosure is observed what is probability that A does not know \(v\)?

\item if non-disclosure is observed what is probability that \(v=1\), that \(v=2\)?

\item what is \(\bar{v}\) when observing non-disclosure?

\item what is A's payoff of not disclosing if \(v=2\)?
\end{itemize}
\end{itemize}
\end{example}
\end{frame}

\begin{frame}[label={sec:orgheadline11}]{Dye's model III}
\begin{itemize}
\item let \(v_i\) be defined such that
\end{itemize}

\(\hspace*{1cm}\mathbb{E}[v|\mbox{ v has no evidence or }v\leq v_i\}] \geq v_i\)
\(\hspace*{1cm}\mathbb{E}[v|\mbox{ v has no evidence or }v\leq v_i\}] \leq v_{i+1}\)

\begin{itemize}
\item this \(v_i\) exists
\begin{itemize}
\item LHS is above RHS for \(v_i = v_1\)
\item LHS is below RHS for \(v_i = v_n\)
\end{itemize}
\end{itemize}

\begin{theorem}[Partial disclosure]
It is an equilibrium that all types above \(v_i\) disclose and those below do not.
\end{theorem}
\end{frame}

\begin{frame}[label={sec:orgheadline12}]{Dye's model IV}
\begin{itemize}
\item say \(v\) is an asset value or stock price that is traded
\item suppose there is a stage 4 where \(v\) becomes known (asset pays out, project finishes, end of year report)
\item how is strategic disclosure different from "disclose everything"?
\begin{itemize}
\item strategic:
\begin{itemize}
\item price of asset drops in stage 2/3 if no disclosure (no news is bad news)
\item without disclosure: uncertainty/volatility about asset value in 4
\end{itemize}
\item disclose everything:
\begin{itemize}
\item price stays constant (in 2/3) if no news
\item price drops if bad news but then no uncertainty about asset value in 4
\end{itemize}
\end{itemize}
\item empirical test for strategic disclosure: does a price drop lead to more volatility?
\end{itemize}
\end{frame}


\section{Limited Communication}
\label{sec:orgheadline16}

\begin{frame}[label={sec:orgheadline14}]{Limited Communication  I}
\begin{itemize}
\item a course has 2 topics
\item student A can either understand or not understand a topic (\(V=\{0,1\}\) for each topic)
\item for each topic the probability of understanding \(f(1)=1/2\) (and values are independent across topics)
\item the exam has limited length: A can communicate at most 1 topic to teacher B
\item with probability \(1-q=1/4\) there is miscommunication and the teacher cannot understand/receive what the student says in the exam (even if A understands the topic)
\item A gets a grade equal to B's ex post expected number of topics A has understood
\end{itemize}
\end{frame}

\begin{frame}[label={sec:orgheadline15}]{Limited Communication  II}
\begin{itemize}
\item what is the probability that "understanding" can be communicated if A chooses the topic of the exam?

\item what is B's belief in case \(v=1\) is (not) communicated if A chooses the topic?

\item what is the probability that "understanding" can be communicated if B chooses the topic of the exam randomly?

\item what is B's belief in case \(v=1\) is (not) communicated if B chooses the topic randomly?

\item does the grade distribution have more variance if A or if B chooses the topic?

\item who should choose the topic if the grades should be as informative as possible about the the student's understanding?
\end{itemize}
\end{frame}

\section{An investment example (BDL)}
\label{sec:orgheadline20}

\begin{frame}[label={sec:orgheadline17}]{Investment decisions I}
\begin{itemize}
\item say A can choose between two technologies
\begin{itemize}
\item technology 1 gives \(v=4\) for sure
\item technology 2 gives \(v=6\) and \(v=0\) with probability 1/2 each
\end{itemize}
\item A's objective is to maximize \(\bar{v}\) (i.e. \(\alpha=0\) for simplicity)
\item A has evidence with probability \(q\) as before
\item the value realizes in stage 4 and goes to B
\end{itemize}

\vspace*{0.5cm}
\begin{itemize}
\item What is the efficient technology choice?
\end{itemize}
\end{frame}

\begin{frame}[label={sec:orgheadline18}]{Investment decisions II}
\begin{itemize}
\item Is it an equilibrium to choose technology 1?
\begin{itemize}
\item in such an equilibrium what is \(\bar{v}\) without disclosure?

\item what is A's payoff when he chooses technology 2?
\end{itemize}

\item Is it an equilibrium to choose technology 2?
\begin{itemize}
\item what will A disclose in such an equilibrium?

\item in such an equilibrium what is \(\bar{v}\) without disclosure?

\item what is A's expected payoff?

\item what is A's expected payoff from deviating to technology 1?
\end{itemize}
\end{itemize}
\end{frame}

\begin{frame}[label={sec:orgheadline19}]{Investment decisions III}
\begin{itemize}
\item general result: A engages in excessive risk taking
\begin{itemize}
\item enjoy upside risk
\item hide downside risk by not disclosing

\item roughly: 
\begin{itemize}
\item \(\hat v\) be the belief if no disclosure
\item A chooses \(\max\{v,\hat v\}\) \(\rightarrow\) convex function of \(v\)
\item A is risk lover
\end{itemize}
\end{itemize}
\end{itemize}
\end{frame}


\section{Revision questions and exercise}
\label{sec:orgheadline24}

\begin{frame}[label={sec:orgheadline21}]{Revision questions}
\begin{itemize}
\item what is information unraveling and why might it occur?
\item what could prevent total information unraveling and why?
\item explain Dye's model setup
\item why is it more informative if the examiner chooses the topic and not the student?
\item why might entrepreneurs choose too risky projects?
\end{itemize}

\vspace*{0.5cm}

\begin{itemize}
\item reading: nothing particularly, I used Eddie Dekel's slides (\url{https://sites.google.com/site/eddiedekelsite/pres_address_2016_slides/pres_address_2016_slides.pdf?attredirects=0&d=1}) and section 5.1.5 in Bolton/Dewatripont's "Contract theory" book as sources. You can find more references in these two sources.
\end{itemize}
\end{frame}

\begin{frame}[label={sec:orgheadline22}]{Exercise I}
An entrepreneur (E) has to sell his project to a perfectly competitive and risk neutral financial market before it matures. With probability \(p\) the project is good and makes a profit of 1. With probability \(1-p\) the project is bad and makes zero profits. With probability \(q\) the entrepreneur finds out whether his project is good before he sells it and in this case he has evidence to document the project value.

\begin{itemize}
\item When will E document the value?

\item What is the market price of the project in case there is (no) evidence?

\item What is the market price of the project in case it is documented to be good?

\item What is the expected market price of the project?
\end{itemize}
\end{frame}

\begin{frame}[label={sec:orgheadline23}]{Exercise II}
\begin{itemize}
\item Say, E is risk averse and obtains Bernoulli utility \(u(m)\) when selling the project at price \(m\) where \(u\) is strictly concave. Does E benefit or suffer from the possibility of learning the project value?

\item Let \(u(m)=\sqrt{m}\). Suppose E has to choose (before the game above is played) one of the following two projects below. Which will he choose? Is this efficient?
\begin{itemize}
\item P1: \(p=1/4\) and \(q=7/8\)
\item P2: \(p=0.16\) and \(q=0\)
\end{itemize}
\end{itemize}
\end{frame}
\end{document}